In general, the approximate factor $q$ of Proposition~11 can be as large as $|Z|$. However, in the following construction, we will construct a matrix where each row has at most $q=3$ distinct values; in this case, Proposition~11 guarantees that $\min_{\mathcal P}\operatorname{Cost}_M(\mathcal P)$ provides a 3-factor approximation of $L_3(M)$.
\begin{theorem}
Given an undirected graph $G=([n],E)$ with $n$ vertices and $|E|=m>0$ edges, one may construct in deterministic polynomial time a function $f_G:[a]\times[n]\to\{0,1,\ldots,\ell\}$, with $\ell=m^2n$, $k=\sqrt{n\ell}=mn$ and $a=\ell+m\cdot k^2$, such that
\[\chi(G)\cdot\ell\le L_3^B(f_G)\le\chi(G)\cdot6\ell.\]
Furthermore, we also have
\[\log\chi(G)+\log\ell\le C_3^B(f_G)\le\log\chi(G)+\log\ell+O(1).\]
\end{theorem}
\begin{proof}
We let $A_\ell$ denote the $\ell\times1$ column vector,
\[A_\ell=\begin{bmatrix}1\\\vdots\\\ell\end{bmatrix}.\]
We let $B_k$ and $C_k$ denote the $k^2\times1$ column vectors
\[B_k=\begin{bmatrix}1\\\vdots\\1\\\vdots\\k\\\vdots\\k\end{bmatrix},\qquad C_k=\begin{bmatrix}1\\\vdots\\k\\\vdots\\1\\\vdots\\k\end{bmatrix},\]
where each value $i\in[k]$ appears $k$ times. Given an edge $\{v,w\}\in[n]\times[n]$, with $v<w$, we define the $k^2\times n$ matrix $M$ so that
\[M_{\{v,w\}}=\begin{bmatrix}0&\cdots&0&B_k&0&\cdots&0&C_k&0&\cdots&0\end{bmatrix},\]
where $B_k$ appears in the $v$-th column and $C_k$ in the $w$-th column. Finally, let $E=\{e_1,\ldots,e_m\}$ denote the edges of $G$; then we define the $a\times n$ communication matrix $f_G$ so that
\[f_G=\begin{bmatrix}A_\ell&\cdots&A_\ell\\\multicolumn{3}{c}{M_{e_1}}\\\multicolumn{3}{c}{\vdots}\\\multicolumn{3}{c}{M_{e_m}}\end{bmatrix}\in\{0,1,\ldots,\ell\}^{a\times n},\]
where $a=\ell+m\cdot k^2$. Observe that each row of $f_G$ has at most $q=3$ distinct values; thus, Proposition~11 provides a 3-factor approximation. We now show the stated inequality, namely, that
\[\chi(G)\cdot\ell\le L_3^B(f_G)\le\chi(G)\cdot2q\ell.\]
The upper-bound is easy to see. Given a coloring of $G$ into $\chi(G)$ colors, we may take the 3-round protocol $\pi$ where Bob first tells Alice which color his vertex has. This partitions the columns by a partition $\mathcal P=\{P_c\}_{c\in[\chi(G)]}$ formed of the various color classes of our coloring. Fix any color $c$ and consider $\operatorname{Cost}_{f_G}(P_c)$. Since $P_c$ is an independent set of $G$, the $C_k$ and $B_k$ columns of each $M_{e_i}$ sub-matrix will always be placed in different parts; therefore, $\operatorname{Cost}_{f_G}(P_c)\le\ell+mk$, where the first term counts the number of rows in $A_\ell$ and the second term counts the number of distinct rows in $M_{e_i}$ restricted to columns of $P_c$ for each $i\in[m]$. Using Proposition~11, we obtain
\[L_3^B(f_G)/q\le\operatorname{Cost}_{f_G}(\mathcal P)=\sum_c\operatorname{Cost}_{f_G}(P_c)\le\chi(G)\cdot(\ell+mk)=\chi(G)\cdot2\ell,\]
and Proposition~11 also gives us
\[C_3^B(f_G)\le\log\chi(G)+\log\ell+O(1).\]
For the other direction, let $\pi$ be any 3-round protocol for $f_G$. The first round of $\pi$ partitions the columns of $f_G$, which is to say, it partitions the vertices of $G$ by some partition $\mathcal P=\{P_i\}_{i\in I}$. We claim that $\chi(G)\cdot\ell\le\operatorname{Cost}_{f_G}(\mathcal P)$ by analyzing the following two cases.
\begin{enumerate}
\item If the partition $\mathcal P$ does not form a coloring of $G$, then there must exist an edge $\{v,w\}$ such that the $v$-th column and the $w$-th column of $G$ are placed in the same part $P_i$. This means that the $C_k$ and $B_k$ columns of the $M_{\{v,w\}}$ sub-matrix are both placed in $P_i$. In this case, we have $\operatorname{Cost}_{f_G}(P_i)\ge k^2\ge n\cdot\ell\ge\chi(G)\cdot\ell$, and thus the lower bound holds.
\item Otherwise, suppose that the partition $\mathcal P$ does form a coloring of $G$. Then the number of parts is $\ge\chi(G)$, and so just the contribution from the first $\ell$ rows gives us $\operatorname{Cost}_{f_G}(\mathcal P)\ge\chi(G)\cdot\ell$.
\end{enumerate}
The lower-bounds on $L_3^B(f_G)$ and $C_3^B(f_G)$ then follow from Proposition~11.
\end{proof}
The following corollary shows that it is NP-hard to approximate $L_3^B(f)$, for a given total two player function $f:[N]\times[N]\to[N]$, with an approximation ratio better than (roughly) $N^{1/5}$.
\begin{corollary}
For every $L\subseteq\{0,1\}^*$ in $\mathrm{NP}$ and every constant $\varepsilon>0$, there exists a polynomial-time algorithm that, on input $x\in\{0,1\}^*$, outputs a communication matrix $M\in[N]^{N\times N}$ such that if $x\in L$ then $L_3^B(M)\le N$, and if $x\notin L$ then $L_3^B(M)>N^{6/5-\varepsilon}$.
\end{corollary}
\begin{proof}
We compose the hardness of approximation result for chromatic number in Theorem~9 with the reduction of Theorem~12, padding the communication matrix with all-0 columns to make it square (since Bob speaks first, this adds at most one leaf), so the communication matrix of $f$ is an $N\times N$ matrix with $N=a=\Theta(n^5)$. The parameters then come from the fact that $n=\Theta(N^{1/5})$.
\end{proof}
Using the bounds on $C_3^B$ instead of the bounds on $L_3^B$ from Theorem~12, we conclude that it is NP-hard to approximate $C_3^B(f)$, for a given total two player function $f:[N]\times[N]\to[N]$, with an approximation ratio better than (roughly) $6/5$. More precisely:
\begin{corollary}
For every $L\subseteq\{0,1\}^*$ in $\mathrm{NP}$ and every $\varepsilon>0$, there exists a polynomial-time algorithm that, on input $x\in\{0,1\}^*$, outputs a communication matrix $M\in[N]^{N\times N}$ such that if $x\in L$ then $C_3^B(M)\le\log N$, and if $x\notin L$ then $C_3^B(M)>(\frac65-\varepsilon)\log N$.
\end{corollary}
